\section{The resource-normalized measurement protocol}
\label{sec:rnm}

\subsection{The contention confound}

Suppose $\n$ validator processes are co-located on a host offering $Q$ CPU cores
in total. If the scheduler divides the host fairly, each validator receives
$Q/\n$ cores, so the per-validator compute budget shrinks as $\n$ grows.
Measured throughput then falls for two unrelated reasons: the growth of
consensus communication, and the shrinking compute budget. An exponent fitted
to such data is a property of the host, not of the protocol.

\begin{Def}[Resource-normalized measurement]
\label{def:rnm}
A measurement campaign satisfies the \emph{\RNM{} protocol} with quota~$\q$ if
every validator container is granted the same CPU quota~$\q$ (in cores) and the
same memory limit, independently of~$\n$, and if the aggregate reservation
$\n\q$ does not exceed the host capacity available to the validator set. The
\emph{normalized throughput} is
\begin{equation}
\TPSn(\n,\q)\eqdef\frac{\TPS(\n,\q)}{\q}.
\label{eq:normalized-tps}
\end{equation}
\end{Def}

Definition~\ref{def:rnm} is implemented with control-group quotas
(\texttt{cpus} and \texttt{mem\_limit} in a container runtime), so it requires
no modification of the consensus software.

\subsection{Invariance property}

\begin{Assumption}[Separable compute budget]
\label{as:separable}
On the operating range considered, throughput factorizes as
\begin{equation}
\TPS(\n,\cc,\q)=\q\cdot\Phi(\n,\cc)\cdot\bigl(1+O(\q)\bigr),
\label{eq:separable}
\end{equation}
that is, the per-validator compute budget enters multiplicatively and does not
interact with the validator count to first order.
\end{Assumption}

Assumption~\ref{as:separable} deserves more than a passing sentence, since
everything downstream rests on it.

\paragraph{Why it should hold.}
A validator's work per block splits into a part proportional to the number of
messages it must verify and a part proportional to the number of transactions
it must execute. Both are CPU-bound and both are throttled by the same cgroup
quota. Under a CFS quota the container receives $\q$ core-seconds per wall
second regardless of how many sibling containers exist, so halving $\q$ halves
the rate at which either part is retired. The validator count enters through
\emph{how much} work there is, the quota through \emph{how fast} it is retired.
These are different arguments of the same production function, which is what
\eqref{eq:separable} states.

\paragraph{Where it should fail.}
The factorization is not universal, and we can name the regimes that break it.
\begin{enumerate*}
\item \emph{Fixed per-process overhead.} Garbage collection, JVM safepoints and
kernel bookkeeping consume a share of $\q$ that does not shrink with $\q$. As
$\q\to0$ this share dominates and throughput falls faster than linearly. The
residual term $O(\q)$ in \eqref{eq:separable} is exactly this effect, and it is
only negligible above some quota floor.
\item \emph{Non-partitioned resources.} A CPU quota does not partition memory
bandwidth, last-level cache or the network stack. Once $\n\q$ approaches the
host capacity, contention reappears through these channels even though the CPU
accounting is clean. Definition~\ref{def:rnm} therefore requires headroom, not
merely $\n\q\le Q$.
\item \emph{Latency-bound operation.} If block production is paced by a fixed
block period rather than by compute, throughput becomes insensitive to $\q$ and
the factorization degenerates. This is a real regime for QBFT at small $\n$ and
is one reason we sweep $\n$ well past the point where the block period stops
binding.
\end{enumerate*}

\paragraph{Status.}
We treat \eqref{eq:separable} as a working hypothesis with a stated domain of
validity, not as a law. The point of the next subsection is that it makes a
prediction sharp enough to be rejected by our own data.

\begin{State}[$\q$-invariance of the consensus exponent]
\label{st:qinvariance}
Under Assumption~\ref{as:separable} and the model
$\Phi(\n,\cc)=\Tzero\,\gfun(\cc)\,\n^{-\beta}$, the least-squares estimate of
$\beta$ obtained by regressing $\log\TPSn$ on $\log\n$ at fixed~$\cc$ does not
depend on the quota~$\q$ used in the campaign.
\end{State}

\begin{Proof}
By~\eqref{eq:normalized-tps} and~\eqref{eq:separable},
$\log\TPSn(\n,\q)=\log\Phi(\n,\cc)+O(\q)$. The term $O(\q)$ is constant across
the regression in~$\log\n$ at fixed~$\q$ and is therefore absorbed by the
intercept. The slope in $\log\n$ equals $-\beta$ for every~$\q$.
\end{Proof}

\begin{Remark}
Statement~\ref{st:qinvariance} converts the quota into an experimental control
variable: choosing a small~$\q$ lowers absolute throughput but leaves the
quantity of interest unchanged, which is what allows a validator set far larger
than the number of physical cores to be studied on one machine.
\end{Remark}

\subsection{Testing the assumption}

Statement~\ref{st:qinvariance} yields a directly testable prediction. Running
the campaign at quotas $\q_{1}<\q_{2}<\q_{3}$ and fitting a model with an
interaction term
\begin{equation}
\log\TPSn = a_{0}+a_{1}\log\n+\sum_{j\ge2}\bigl(b_{j}+d_{j}\log\n\bigr)
\mathbf{1}\{\q=\q_{j}\}+\epsilon
\label{eq:ancova}
\end{equation}
the invariance hypothesis is $H_{0}: d_{2}=d_{3}=0$, tested by a partial
$F$-test. Rejection of $H_{0}$ indicates that
Assumption~\ref{as:separable} fails and that the measured exponent is
quota-dependent.

A second and more direct test targets the factorization itself rather than its
consequence for the slope. Under \eqref{eq:separable}, raw throughput at fixed
$(\n,\cc)$ should be proportional to $\q$ with a zero intercept; a significantly
positive intercept in the regression of $\TPS$ on $\q$ indicates a
latency-bound regime, and a significantly negative one indicates fixed overhead
eating into the quota. We report both tests in Section~\ref{sec:results}.

If either test rejects, the honest response is not to discard the protocol but
to restrict it: the invariance claim then holds only on the quota interval where
it survives, and the exponent must be reported together with that interval.
